\makeabstract
{
Abundance of BK Channels in Mouse Inner Hair Cells Throughout the Onset of Hearing
}
{
Hannah Feng (1), Juliana Giacomini (2), Radha Kalluri (2)
}
{
\textsuperscript{1}Amherst College, Amherst, MA, USA\\
\textsuperscript{2}Caruso Department of Otolaryngology-Head and Neck Surgery, Keck School of Medicine of University of Southern California, CA, USA
}
{
Inner hair cells (IHCs) are the primary sensory receptors for hearing in the cochlea of mammals, and their expression of ion channels changes meticulously throughout hearing development. Among these channels, large-conductance calcium- and voltage-gated potassium (BK) channels are of particular interest as they generate the electrical currents that are necessary for mature hearing. This study aims to quantify the abundance of BK channels in mouse IHCs before, during, and after the onset of hearing in order to better understand their developmental role.
}
{
Because mice are born deaf and acquire hearing around postnatal day (P)10-12, three cohorts of mice (ages P8-P14) were euthanized for cochlear tissue collection. Western blotting was performed on nine samples to detect BK and glyceraldehyde 3-phosphate dehydrogenase (GAPDH) as a loading control. Band intensities were measured in ImageJ, corrected for background, and normalized to GAPDH. A standard curve was generated by loading increasing amounts of BK protein, confirming a linear relationship between band intensity and protein load. Antibody specificity was verified using BK knockout tissue.
}
{
BK protein was detectable in three of nine samples, reflecting variability in tissue dissection and preparation. Expression was observed as early as P8, prior to hearing onset. In this sample, normalized BK intensity rose from 0.028 (P8) to 0.093 (P11; +232\%). In two additional samples, BK levels increased from 0.067 (P10) to 0.212 (P14; +216\%) and from 0.0295 (P10) to 0.0903 (P13; +206\%). Due to small sample size and high variability between replicates, regression analysis was not feasible.
}
{
BK channels are detectable in mouse IHCs before the onset of hearing, suggesting a previously unrecognized role in development. Expression appears to increase with maturation, indicating potential involvement in the transition to functional hearing. These findings underscore the need for further study of ion channel regulation in IHC development, as such insights may inform gene therapy strategies aimed at regenerating hair cells and restoring hearing function.
}

\newpage
